\chapter{Oprogramowanie}\label{ch:oprogramowanie}

\begin{sekcja}{Narzędzia}
\renewcommand{\arraystretch}{1.3}
\begin{tabularx}{\linewidth}{@{}P{34mm}LP{50mm}@{}}
\toprule
\textbf{Narzędzie} & \textbf{Zastosowanie} & \textbf{Adres}\\\midrule
SondeHub Burst Calculator & Ilość gazu, siła nośna, wysokość pęknięcia & \www{sondehub.org/calc}\\
SondeHub Predictor & Predykcja trajektorii i miejsca lądowania & \www{predict.sondehub.org}\\
SondeHub Amateur & Mapa śledzenia (APRS, RTTY, LoRa, Horus) & \www{amateur.sondehub.org}\\
aprs.fi & Mapa śledzenia APRS & \www{aprs.fi}\\
PureTrack & Mapa śledzenia & \www{puretrack.io}\\
SPOT & Lokalizator satelitarny & \www{www.findmespot.com}\\
Find My / Find Hub & Znaczniki Bluetooth & aplikacje Apple i Google\\
\bottomrule
\end{tabularx}
\end{sekcja}

\begin{sekcja}{Konfiguracja przed misją}
Wykonaj najpóźniej \textbf{2 dni przed lotem} (formularz F7):
\begin{enumerate}
  \item Wejdź na \www{amateur.sondehub.org} i skonfiguruj dokument ładunku (\emph{payload}) dla RTTY. Bez tego SondeHub nie zdekoduje telemetrii.
  \item Włącz tracker na zewnątrz. Wejdź na \www{aprs.fi} i sprawdź, czy pozycja znaku z SSID pojawia się na mapie.
  \item Wejdź na SondeHub Amateur i sprawdź, czy pakiety RTTY pojawiają się na mapie.
  \item Aktywuj subskrypcję SPOT (plan Flex -- na jeden miesiąc). Wejdź na stronę SPOT i sprawdź pozycję urządzenia.
  \item Dodaj znaczniki Bluetooth do kont Apple i Google zespołu. Sprawdź ich pozycję w aplikacjach.
  \item Udostępnij linki do śledzenia całemu zespołowi.
\end{enumerate}
\end{sekcja}

\begin{sekcja}{Wersja kodu lotnego}
Kod lotny (HABduino, OBC, ładunek) i dokumentacja tworzą jedną całość:
\begin{itemize}
  \item w dokumentacji misji (karta misji F0, raport) zapisz wersję kodu lotnego: numer wydania (\emph{release}) albo commita,
  \item do wydania kodu lotnego dołącz końcową wersję dokumentacji (np. PDF tego podręcznika i raport z misji),
  \item nie zmieniaj kodu po przeglądzie gotowości. Każda zmiana po przeglądzie wymaga ponownych testów.
\end{itemize}

Gdy po przeglądzie znajdziesz błąd:
\begin{itemize}
  \item \textbf{błąd krytyczny} (zagraża bezpieczeństwu albo celowi misji): popraw kod, powtórz testy i przegląd gotowości albo przełóż lot,
  \item \textbf{błąd niekrytyczny}: nie zmieniaj kodu. Bardziej wartościowe jest zrozumienie, kiedy błąd występuje, i opracowanie procedury obejścia poza kodem -- np. zmiana sekwencji uruchomienia albo kolejności telekomend. Zapisz obejście w formularzach dnia startu i w raporcie.
\end{itemize}
\end{sekcja}

\begin{sekcja}{Konfiguracja komputerów bez przeprogramowania}\label{sec:konfiguracja}
Kod lotny nie zmienia się po przeglądzie gotowości, ale parametry misji tak: wysokość i czas zadziałania odcinacza, granice dozwolonego obszaru, parametry lotu, godzina startu. Zespół musi wpisywać te wartości \textbf{bez przeprogramowania} komputera.

Zespół tworzy własną procedurę obsługi każdego komputera (HABduino, OBC, ładunek). Procedura określa, jak wpisać wartości i kto to robi. Możliwe sposoby:
\begin{itemize}
  \item kabel i komendy, które zapisują konfigurację w pamięci trwałej komputera,
  \item telekomendy -- tylko gdy komputer ma odbiornik,
  \item karta SD z plikiem konfiguracyjnym, wymieniana przed startem.
\end{itemize}

\textbf{Zasady:}
\begin{enumerate}
  \item Przygotuj konfigurację dla \textbf{każdego scenariusza misji} (np. miejsce startu zasadnicze i zapasowe, opóźniony start). Nadaj każdej konfiguracji nazwę scenariusza i zapisz ją w dokumentacji misji.
  \item Przetestuj każdą konfigurację przed lotem, tą samą drogą, którą wpiszesz ją w dniu startu.
  \item Po wpisaniu odczytaj konfigurację z komputera i porównaj z zapisaną. Nie sprawdzaj pliku na laptopie -- sprawdź to, co jest w komputerze.
  \item Ustawień, które zmieniają się w locie (z wysokością albo czasem), test na ziemi nie sprawdzi. Przeczytaj je w konfiguracji (sekcja~\ref{sec:konfiguracja-trackera}).
  \item Zapisz w formularzach dnia startu, którą konfigurację wpisano do którego komputera. Zachowaj pliki konfiguracyjne z lotu do późniejszej weryfikacji.
\end{enumerate}
\end{sekcja}
